
\section{The Generalized Road Coloring Problem}
The definitions below allow a general constant outdegree $d$ so that the
historical formulations can be read precisely. In our road-coloring
applications we specialize to the simple, loopless, irreducible and
aperiodic two-out setting stated at the opening. The generalized periodic
formulation at the end is recorded separately and does not relax those
application hypotheses.

\label{sec:1.1}

Let $\mathcal A$ be the adjacency matrix of a strongly connected digraph $G=(V,\E)$.  
We will assume that the row sums of $\mathcal A$ are constant or, equivalently, that 
the outdegree of each vertex of $G$ is constant.

\begin{definition}	\label{defn:1}
% BEGIN SUBJECT INDEX: defn:1
\index{coloring!definition}
\index{transformation!matrix representation}
\index{right action}
% END SUBJECT INDEX

A coloring of $G$ is a decomposition of $\mathcal A$ into $d$ binary (that is, 0--1) 
stochastic matrices $R_1,R_2, ..., R_d$ such that $\mathcal A = R_1 + .... + R_d$.

Intuitively, one can think of  a matrix $R_i$ as assigning the $i$th of $d$ 
colors to the edge $(j,k)$ if $(R_i)_{j,k} = 1$. 
Each $n \times n$ binary stochastic matrix $R$ can also be thought of as a function on the vertices $\{1,2,...,n\}$ where $R_{j,k} = 1$ iff $jR = k$. (We will write functions to the right of domain elements to duplicate the order of matrix multiplication.)
	
Let $S = \langle R_1, R_2, ..., R_d \rangle$ be the semigroup generated by 
$R_1, R_2, ...,R_d$ under multiplication.  
Clearly, $S$ is a finite semigroup.  
We will refer to $S$ as a coloring semigroup, and  will write $E(S)$ for the 
idempotents in $S$.  
\end{definition}

\begin{definition}[Target-pair graph]
\label{def:target-pair-graph}
% BEGIN SUBJECT INDEX: def:target-pair-graph
\index{graph!target-pair}
\index{target-pair graph}
% END SUBJECT INDEX

For a two-out coloring $(R,B)$, use $A=(R+B)/2$. On
$V=\{1,\ldots,n\}$, put
\[
 t_x=\{xR,xB\},\qquad
 \mathcal T(A)=(t_1,\ldots,t_n),\qquad
 H_A=(V,\{t_x:x\in V\}).
\]
Here $H_A$ is undirected: each source contributes an edge joining its two
\emph{destinations}, not an edge from the source to a destination.
The source-indexed list $\mathcal T(A)$ retains repeated auxiliary edges;
$H_A$ suppresses their multiplicities when discussing connectivity or
bipartiteness. For single-digit vertex labels, write $ij$ for the unordered
pair $\{i,j\}$, with $i<j$, to display the list compactly.
\end{definition}

\begin{remark}
The target-pair graph is invariant under every interchange of the two
outgoing colors at a source; it depends on the uncolored two-out graph,
not its coloring. Under the simple road-coloring hypotheses it is loopless,
although different sources may yield the same auxiliary edge. In the more
general two-generator setting, a coincident pair $xR=xB$ contributes a loop.
The connected components, their bipartiteness, and hence the ranks of $A$
and $\Delta$ are fixed throughout a recoloring run. The two bipartition
sides, when they exist, describe CC target types; they are not generally
kernel blocks. The component partition is the graph-intrinsic canonical
partition. These distinctions will be used in the rank and CC discussions.
\end{remark}



\begin{definition}	\label{defn:2}
	Let $S = \langle R_1,R_2,\ldots,R_d\rangle$ be a coloring semigroup for the digraph $G$.
We will say that $S$ has a synchronizing instruction, if $S$ contains a matrix of rank one.
\end{definition}

Notice that a rank one binary stochastic matrix has a column $k$ of all ones.  
We can think of this matrix as a function mapping every vertex to vertex $k$, which intuitively captures the idea of synchronization.

The kernel $K$ of a finite transformation semigroup consists of its
minimum-rank transformations. Its rank is one precisely when the coloring
is synchronizing. The structure of completely simple kernels, their Rees
coordinates, partitions, ranges, local groups and right groups will be
developed systematically in the semigroup chapter; we do not need their
full structure merely to state the coloring questions.

The final concept required to precisely state the road coloring problem is the notion of the period of a graph.

\begin{definition}    \label{defn:3}  
	Let $\mathcal P = \{ P_1,P_2,\ldots,P_t \}$ be a partition of the vertices $V$ such that if $(i,j) \in \E$ and $i \in P_k$, then $j \in P_{k+1}$ ($i \in P_t$ implies $j \in P_1$).  Then $\mathcal P$ will be called a periodic partition of order $t \ge 2$.
\end{definition}

\begin{definition}    \label{defn:4}  
	A digraph $G=(V,\E)$ is periodic of period $t\ge 2$ iff 
$$
	t = \max_\mathcal P [\hbox{card}(\mathcal P)]
$$
where this maximum is taken over all periodic partitions of $G$. 
If $G$ has no periodic partitions, then $G$ is aperiodic.
In this case, we will refer to $\mathcal P = \{V\}$ as the periodic partition.
\end{definition}

The above definition can be shown to be equivalent to the more familiar one that uses the period of a vertex, where 
per$(v) = \gcd\{|c_1|,\ldots,|c_n|\}$ where $c_1,\ldots,c_n$ are the simple cycles containing $v$.

	The road coloring problem can now be stated in the following way.

\runinhead{Classical road coloring statement.}
	Let $G=(V,\mathcal E)$ be a $d$-out, strongly connected digraph 
with adjacency matrix $\mathcal A$. Let $\left\{\left\{R_{i_1},R_{i_2},\dots,R_{i_d}\right\}:
\mathcal A=R_{i_1}+R_{i_2}+ \ldots +R_{i_d}\right\}$ be the set of all colorings of 
$G$ and let $S_i=\left\langle R_{i_1},R_{i_2},\dots,R_{i_d}\right\rangle$ 
be the  associated coloring semigroups. If $G$ is aperiodic, then for some 
$i$, $S_i$ contains a synchronizing instruction.


	In order to generalize the road coloring conjecture to all strongly 
connected graphs, one needs to observe that if the completely simple kernel 
of a semigroup of functions is a right group, then each element in the kernel 
induces the same partition on the set of vertices.  Here, of course, the partition 
$\mathcal P$ induced by a function $k$ is the collection of subsets of $V$ of the form 
$k^{-1}(v)$, for $v \in V$.  In the case of the constant functions, the partition 
induced by the kernel is simply $\{V\}$.

\runinhead{Generalized Road Coloring Conjecture.}
	Let $G=(V,\mathcal E)$ be a strongly connected $d$-out digraph. Then $G$ is
periodic of period $t\ge 1$ with maximal periodic partition 
$\mathcal P=\left\{P_1,\dots,P_t\right\}$ iff $G$ has a minimal 
{\rm(}with respect to the rank of its kernel{\rm )}
coloring semigroup $S=\left\langle R_1,\dots,R_d\right\rangle$ whose kernel
$K$ is a right group with partition $\mathcal P$.


	Clearly if $\mathcal P = \{P_1,\ldots,P_t\}$ is the partition induced by the kernel $K$, then rank$(K)=t$. The semigroup discussion explains why having a minimal coloring 
semigroup whose kernel is a right group with rank $t$ is a sufficient condition for the graph to be periodic of order $t$. 


